% Part: computability
% Chapter: recursive-functions
% Section: sequences

\documentclass[../../../include/open-logic-section]{subfiles}

\begin{document}

\olfileid[pt]{cmp}{rec}{seq}
\olsection{Sequências}

O conjunto das funções recursivas primitivas é notavelmente robusto.
Mas poderemos fazer ainda mais depois de desenvolvermos um meio adequado
de lidar com \emph{sequências}. Identificaremos sequências finitas de
números naturais com números naturais da seguinte maneira: a sequência
$\langle a_0, a_1, a_2, \dots, a_k \rangle$ corresponde ao número
\[
p_0^{a_0+1} \cdot p_1^{a_1+1} \cdot p_2^{a_2+1} \cdot \dots \cdot
p_k^{a_k+1}.
\]
Somamos um aos expoentes para garantir que, por exemplo, as sequências
$\langle 2, 7, 3\rangle$ e $\langle 2, 7, 3, 0, 0 \rangle$ tenham
códigos numéricos distintos. Podemos usar tanto $0$ quanto~$1$ para
codificar a sequência vazia; para sermos específicos, seja $\emptyseq$
uma notação para~$0$.

A razão pela qual essa codificação de sequências funciona é o chamado
Teorema Fundamental da Aritmética: todo número natural $n \ge 2$ pode
ser escrito de uma, e apenas uma, maneira na forma
\[
n = p_0^{a_0} \cdot p_1^{a_1} \cdot \dots \cdot p_k^{a_k}
\]
com $a_k \ge 1$. Isso garante que a aplicação $\tuple{}(a_0,
\dots, a_k) = \tuple{a_0, \dots, a_k}$ é injetora: sequências
diferentes são associadas a números diferentes; a cada número
corresponde, no máximo, uma sequência.

Mostraremos agora que as operações de determinar o comprimento de uma
sequência, determinar seu $i$-ésimo elemento, acrescentar um elemento a
uma sequência e concatenar duas sequências são todas recursivas primitivas.

\begin{prop}
  A função $\len{s}$, que devolve o comprimento da sequência
  $s$, é recursiva primitiva.
\end{prop}

\begin{proof}
  Seja $R(i, s)$ a relação definida por
  \[
  R(i, s) \text{ se, e somente se, }
  p_i \mid s \land p_{i+1} \nmid s.
  \]
  $R$ é claramente recursiva primitiva. Sempre que $s$ é o código de
  uma sequência não vazia, O.s.\@,
  \[
  s = p_0^{a_0+1} \cdot \dots \cdot p_{k}^{a_{k}+1},
  \]
  $R(i,s)$ vale se $p_i$ é o maior primo tal que
  $p_i \mid s$, O.s.\@, $i = k$. Assim, o comprimento de $s$ é $i+1$
  se, e somente se, $p_i$ é o maior primo que divide~$s$; portanto,
  podemos definir
  \[
  \len{s} =
  \begin{cases}
    0 & \text{se $s = 0$ ou $s = 1$} \\
    1 + \bmin{i < s}{R(i, s)} & \text{caso contrário}
  \end{cases}
  \]
  Podemos usar aqui a minimização limitada, pois há apenas um $i$ que
  satisfaz $R(i,s)$ quando $s$~é o código de uma sequência e, se $i$
  existe, ele é menor do que o próprio~$s$.
\end{proof}

\begin{prop}
  A função $\fn{append}(s,a)$, que devolve o resultado de acrescentar $a$ à
  sequência $s$, é recursiva primitiva.
\end{prop}

\begin{proof}
  $\fn{append}$ pode ser definida por:
  \[
  \fn{append}(s,a) =
  \begin{cases}
    2^{a+1} & \text{se $s = 0$ ou $s = 1$} \\
    s \cdot p_{\len{s}}^{a+1} & \text{caso contrário.}
  \end{cases}
  \]
\end{proof}

\begin{prop}
  A função $\fn{element}(s,i)$, que devolve o $i$-ésimo elemento de $s$
  (sendo o elemento inicial chamado de $0$-ésimo), ou $0$ se $i$ é
  maior ou igual ao comprimento de $s$, é recursiva primitiva.
\end{prop}

\begin{proof}
Observe que $a$ é o $i$-ésimo elemento de~$s$ se, e somente se,
$p_i^{a+1}$ é a maior potência de~$p_i$ que divide~$s$, O.s.\@,
$p_i^{a+1} \mid s$ mas $p_i^{a+2} \nmid s$. Portanto:
  \[
  \fn{element}(s,i) =
  \begin{cases}
    0 & \mbox{se $i \geq \len{s}$} \\
    \bmin{a < s}{(p_i^{a+2} \nmid s)} & \text{caso contrário.}
  \end{cases}
  \]
\end{proof}

Em vez de usar os nomes oficiais das funções definidas acima,
introduzimos uma notação mais compacta. Usaremos $(s)_i$ em vez de
$\fn{element}(s,i)$, e $\tuple{s_0, \dots, s_k}$ como abreviação de
\[
\fn{append}(\fn{append}(\dots \fn{append}(\emptyseq,s_0)
\dots),s_k).
\]
Observe que, se $s$ tem comprimento~$k$, os elementos de $s$ são
$(s)_0$, \dots,~$(s)_{k-1}$.

\begin{prop}
A função $\fn{concat}(s,t)$, que concatena duas
sequências, é recursiva primitiva.
\end{prop}

\begin{proof}
  Queremos uma função $\fn{concat}$ com a propriedade de que
  \[
    \fn{concat}(\tuple{a_0, \dots, a_k}, \tuple{b_0, \dots, b_l}) =
    \tuple{a_0, \dots, a_k, b_0, \dots, b_l}.
  \]
  Usaremos uma função ``auxiliar''
  $\fn{hconcat}(s,t,n)$ que concatena os primeiros $n$ símbolos de $t$
  a~$s$. Essa função pode ser definida por recursão primitiva da
  seguinte maneira:
  \begin{align*}
    \fn{hconcat}(s,t,0) & = s\\
    \fn{hconcat}(s,t,n+1) & = \fn{append}(\fn{hconcat}(s,t,n),(t)_n)
    \intertext{Podemos então definir $\fn{concat}$ por}
    \fn{concat}(s,t) & = \fn{hconcat}(s,t,\len{t}).
  \end{align*}
\end{proof}

Escreveremos $s \concat t$ em vez de $\fn{concat}(s,t)$.

Será útil podermos limitar o código numérico de uma sequência em termos
de seu comprimento e de seu maior elemento. Suponha que $s$ seja uma
sequência de comprimento~$k$ cujos elementos sejam todos menores ou
iguais a certo número~$x$. Então $s$ tem, no máximo, $k$ fatores primos,
cada um deles, no máximo,~$p_{k-1}$, e cada um elevado, no máximo, a
$x+1$ na fatoração prima de~$s$. Em outras palavras, se definirmos
\[
\fn{sequenceBound}(x,k) = p_{k-1}^{k \cdot (x+1)},
\]
então o código numérico da sequência~$s$ descrita acima é, no máximo,
\[
\fn{sequenceBound}(x,k).
\]

Dispor de tal limite para sequências nos oferece uma maneira de definir
novas funções por meio de busca limitada. Por exemplo, podemos definir
$\fn{concat}$ usando busca limitada. Tudo o que precisamos fazer é
escrever uma \emph{especificação} recursiva primitiva do objeto (o número
da sequência concatenada) que procuramos, além de um limite até onde
procurar. A seguinte definição funciona:
\begin{align*}
  \fn{concat}(s,t) = {} & \bmin{v < \fn{sequenceBound}(s+t,\len{s} +
    \len{t})}{} \\
  & \quad(\len{v} = \len{s} + \len{t} \land {}\\
    & \qquad \bforall{i < \len{s}}{((v)_i = (s)_i) \land {} \\
      & \qquad \bforall{j < \len{t}}{((v)_{\len{s}+j} = (t)_j)})}
\end{align*}

\begin{prob}
Mostre que existe uma função recursiva primitiva~$\fn{sconcat}(s)$
com a propriedade de que
\[
\fn{sconcat}(\tuple{s_0, \dots, s_k}) = s_0 \concat \dots \concat s_k.
\]
\end{prob}

\begin{prob}
Mostre que existe uma função recursiva primitiva~$\fn{tail}(s)$
com a propriedade de que
\begin{align*}
  \fn{tail}(\emptyseq) & = 0 \text{ e}\\
  \fn{tail}(\tuple{s_0, \dots, s_{k}}) & = \tuple{s_1, \dots, s_{k}}.
\end{align*}
\end{prob}

\begin{prop}
  \ollabel{prop:subseq}
  A função $\fn{subseq}(s, i, n)$, que devolve a subsequência
  de $s$ de comprimento~$n$ iniciada no $i$-ésimo elemento, é recursiva
  primitiva.
\end{prop}

\begin{proof}
  Exercício.
\end{proof}

\begin{prob}
  Demonstre \olref[cmp][rec][seq]{prop:subseq}.
\end{prob}
  
\end{document}
